\section{Run dependence}
\label{sec:rundep}

The period split was motivated by the observation that the fitted curvature term differs between subsets of the sample. Fits with
the KFP mass (no $\kappa$) on run groups (Table~\ref{tab:groups}) and per-run counts (Table~\ref{tab:runs}) show:
\begin{itemize}
\item $a$ grows with the run period: \GrpThirtyOne\%/(GeV/$c$) on run 79509 segments 0--30, \PerAKfpKsA{} on period A,
  \GrpNine{} on runs 80904--80915 and \GrpEightOne{} on runs 81556--81609;
\item the number of \Ks{} QA candidates per event (signal plus combinatorial background) grows in the same order:
  \PerAKsEvMin--\PerAKsEvMax{} per run in period A and \PerBKsEvMin--\PerBKsEvMax{} in period B. This points to higher
  luminosity or pile-up in the later runs; the luminosity itself is not checked here.
\end{itemize}
A luminosity-dependent effect such as TPC space-charge distortions would produce this pattern, but this note does not test the
hypothesis. The $\kappa$ table, derived in period A, removes the same fraction of $a$ in both periods (Sec.~\ref{sec:results}).

\begin{table}[H]
\centering\small
\caption{\Ks{} fit with the KFP mass on run groups (same method; $s$ grid 0.975--1.015). Statistical intervals only. Logs
\fpath{logs/fitMomRes_{kfp31,g7951x,g809xx,g81xxx,k287_*_kfp}.txt}.}
\label{tab:groups}
\begin{tabular}{lcccc}
\toprule
runs & $a$ (\%/(GeV/$c$)) & $b$ (\%) & $s$ & $\chi^2$/ndf \\
\midrule
% generated by scripts/make_numbers.py
79509 seg.\ 0--30 & $2.00^{+0.76}_{-0.26}$ & $1.80^{+0.07}_{-0.24}$ & 0.990 & 427.2/389 \\
79509--79516 (A) & $3.00^{+0.06}_{-0.27}$ & $1.60^{+0.03}_{-0.05}$ & 0.991 & 470.4/389 \\
79510--79516 & $3.00^{+0.13}_{-0.15}$ & $1.60^{+0.02}_{-0.16}$ & 0.991 & 510.2/389 \\
80904--80915 & $3.75^{+0.18}_{-0.34}$ & $1.20^{+0.10}_{-0.10}$ & 0.991 & 395.6/389 \\
81556--81609 & $4.25^{+0.79}_{-0.31}$ & $1.40^{+0.10}_{-0.48}$ & 0.990 & 427.4/389 \\
B (all non-A) & $4.00^{+0.29}_{-0.18}$ & $1.20^{+0.15}_{-0.05}$ & 0.991 & 400.5/389 \\
all & $3.75^{+0.07}_{-0.27}$ & $1.40^{+0.05}_{-0.05}$ & 0.991 & 421.5/389 \\

\bottomrule
\end{tabular}
\end{table}

\begin{table}[H]
\centering\scriptsize
\caption{Per run: period, segments, processed events, \Ks{} QA candidates (all masses in 0.40--0.60~GeV/$c^2$, signal + background)
and per event, \La+$\bar\Lambda$ QA candidates, and mean $\kappa$ of the QA candidates (\Ks{} $\pi^-$, $\pi^+$; \La{} pion,
proton). From the condor job logs (\fpath{logs/joblines.txt}).}
\label{tab:runs}
\setlength{\tabcolsep}{3.5pt}
\begin{tabular}{rcrrrcrcccc}
\toprule
run & per. & seg. & events & \Ks{} cand. & \Ks/ev & \La{} cand. & $\kappa_{\pi^-}$ & $\kappa_{\pi^+}$ & $\kappa_{\pi}^{\Lambda}$ & $\kappa_{p}^{\Lambda}$ \\
\midrule
% generated by scripts/make_numbers.py
79509 & A & 39 & 38{,}656 & 8{,}939 & 0.231 & 17{,}656 & 0.9980 & 1.0164 & 1.0078 & 1.0067 \\
79510 & A & 8 & 7{,}931 & 1{,}599 & 0.202 & 3{,}306 & 0.9978 & 1.0160 & 1.0077 & 1.0067 \\
79511 & A & 8 & 7{,}934 & 1{,}717 & 0.216 & 3{,}324 & 0.9980 & 1.0160 & 1.0076 & 1.0063 \\
79512 & A & 8 & 7{,}932 & 1{,}663 & 0.210 & 3{,}087 & 0.9969 & 1.0167 & 1.0076 & 1.0063 \\
79513 & A & 8 & 7{,}931 & 1{,}619 & 0.204 & 2{,}942 & 0.9975 & 1.0164 & 1.0079 & 1.0063 \\
79514 & A & 8 & 7{,}931 & 1{,}493 & 0.188 & 2{,}783 & 0.9977 & 1.0161 & 1.0076 & 1.0063 \\
79515 & A & 8 & 7{,}930 & 1{,}473 & 0.186 & 2{,}623 & 0.9974 & 1.0160 & 1.0079 & 1.0052 \\
79516 & A & 8 & 7{,}928 & 1{,}338 & 0.169 & 2{,}379 & 0.9987 & 1.0162 & 1.0076 & 1.0072 \\
80904 & B & 8 & 7{,}936 & 3{,}778 & 0.476 & 8{,}545 & 0.9971 & 1.0160 & 1.0076 & 1.0063 \\
80905 & B & 8 & 7{,}938 & 3{,}572 & 0.450 & 8{,}029 & 0.9974 & 1.0162 & 1.0078 & 1.0064 \\
80906 & B & 8 & 7{,}938 & 3{,}601 & 0.454 & 7{,}882 & 0.9974 & 1.0160 & 1.0077 & 1.0066 \\
80908 & B & 8 & 7{,}939 & 3{,}036 & 0.382 & 6{,}692 & 0.9983 & 1.0158 & 1.0075 & 1.0066 \\
80909 & B & 8 & 7{,}936 & 2{,}666 & 0.336 & 6{,}117 & 0.9968 & 1.0162 & 1.0077 & 1.0058 \\
80913 & B & 8 & 7{,}930 & 2{,}166 & 0.273 & 4{,}397 & 0.9971 & 1.0162 & 1.0077 & 1.0060 \\
80915 & B & 8 & 7{,}933 & 1{,}994 & 0.251 & 3{,}796 & 0.9972 & 1.0161 & 1.0073 & 1.0064 \\
81556 & B & 8 & 7{,}936 & 7{,}515 & 0.947 & 17{,}241 & 0.9978 & 1.0163 & 1.0076 & 1.0068 \\
81558 & B & 8 & 7{,}937 & 7{,}086 & 0.893 & 16{,}020 & 0.9981 & 1.0163 & 1.0077 & 1.0070 \\
81560 & B & 8 & 7{,}941 & 4{,}724 & 0.595 & 10{,}797 & 0.9980 & 1.0157 & 1.0077 & 1.0068 \\
81562 & B & 8 & 7{,}939 & 4{,}152 & 0.523 & 9{,}505 & 0.9980 & 1.0162 & 1.0077 & 1.0067 \\
81566 & B & 8 & 7{,}933 & 3{,}116 & 0.393 & 6{,}854 & 0.9975 & 1.0159 & 1.0078 & 1.0062 \\
81575 & B & 8 & 7{,}938 & 6{,}779 & 0.854 & 16{,}118 & 0.9978 & 1.0165 & 1.0079 & 1.0069 \\
81576 & B & 8 & 7{,}936 & 6{,}319 & 0.796 & 14{,}385 & 0.9977 & 1.0162 & 1.0077 & 1.0068 \\
81577 & B & 8 & 7{,}937 & 6{,}031 & 0.760 & 13{,}923 & 0.9979 & 1.0165 & 1.0077 & 1.0069 \\
81580 & B & 8 & 7{,}940 & 4{,}564 & 0.575 & 10{,}359 & 0.9978 & 1.0156 & 1.0076 & 1.0066 \\
81581 & B & 8 & 7{,}940 & 4{,}384 & 0.552 & 9{,}782 & 0.9979 & 1.0165 & 1.0077 & 1.0069 \\
81583 & B & 8 & 7{,}937 & 3{,}930 & 0.495 & 8{,}828 & 0.9977 & 1.0165 & 1.0079 & 1.0064 \\
81584 & B & 8 & 7{,}939 & 3{,}446 & 0.434 & 7{,}746 & 0.9978 & 1.0160 & 1.0078 & 1.0061 \\
81585 & B & 8 & 7{,}933 & 3{,}335 & 0.420 & 7{,}403 & 0.9983 & 1.0159 & 1.0078 & 1.0067 \\
81586 & B & 8 & 7{,}931 & 3{,}403 & 0.429 & 7{,}605 & 0.9981 & 1.0159 & 1.0079 & 1.0065 \\
81600 & B & 8 & 7{,}935 & 7{,}773 & 0.980 & 18{,}212 & 0.9974 & 1.0163 & 1.0078 & 1.0066 \\
81601 & B & 8 & 7{,}937 & 6{,}193 & 0.780 & 14{,}406 & 0.9977 & 1.0161 & 1.0076 & 1.0067 \\
81609 & B & 8 & 7{,}936 & 3{,}855 & 0.486 & 8{,}527 & 0.9972 & 1.0160 & 1.0078 & 1.0064 \\

\bottomrule
\end{tabular}
\end{table}
